% Fichier généré automatiquement par tools/generate_corrections.py.
% Source : DYN/DYN-04-TorseurDynamique/STOCK/03_TT_02/03_TT_02.tex
% Statut : complet (3/3 question(s) renseignée(s), 3 reprise(s)).
\begin{corrigebox}[Corrigé extrait de la banque]
\CorrigeQuestion{1}
\CorrigeSource{Réponse reprise d'un exercice homonyme de la banque ; la provenance exacte est indiquée dans le commentaire du fichier.}
% Réponse reprise de : CIN/CIN-02-VitesseAcceleration/03_TT_02/03_TT_02.tex
Par dérivation vectorielle, on a : $\vectv{C}{2}{0} $
$ = \deriv{\vect{AC}}{\rep{0}}$
$ = \dot{\lambda}(t)\vect{i_0}+\dot{\mu}(t)\vect{j_0}$.

Par composition du torseur cinématique, on a : 
$\vectv{C}{2}{0} = \vectv{C}{2}{1} +\vectv{C}{1}{0}$
$ = \deriv{\vect{BC}}{\rep{1}}+\deriv{\vect{AC}}{\rep{0}}$
$ = \dot{\lambda}(t)\vect{i_0}+\dot{\mu}(t)\vect{j_0}$.
\CorrigeQuestion{2}
\CorrigeSource{Réponse reprise d'un exercice homonyme de la banque ; la provenance exacte est indiquée dans le commentaire du fichier.}
% Réponse reprise de : CIN/CIN-02-VitesseAcceleration/03_TT_02/03_TT_02.tex
$\torseurcin{V}{2}{0}= \torseurl{\vect{0}}{\dot{\lambda}(t)\vect{i_0}+\dot{\mu}(t)\vect{j_0}}{\forall P}$.
\CorrigeQuestion{3}
\CorrigeSource{Réponse reprise d'un exercice homonyme de la banque ; la provenance exacte est indiquée dans le commentaire du fichier.}
% Réponse reprise de : CIN/CIN-02-VitesseAcceleration/03_TT_02/03_TT_02.tex
$\vectg{C}{2}{0} = \deriv{\vectv{C}{2}{0}}{\rep{0}}=\ddot{\lambda}(t)\vect{i_0}+\ddot{\mu}(t)\vect{j_0}$.
\end{corrigebox}
